\section{Introducción}

\section{Introducción}

\subsection{Objetivos del Trabajo Práctico}

El objetivo de este trabajo práctico es diseñar, implementar y validar un procesador \textit{RISC-V} segmentado de cinco etapas sobre FPGA, incorporando además una interfaz de depuración y control externo mediante comunicación UART. 

En primer lugar, se pretende construir un \textit{pipeline} funcional capaz de ejecutar un subconjunto de instrucciones de la arquitectura RISC-V, asegurando el correcto transporte de datos y señales de control entre las distintas etapas. Esto implica no solo el diseño del camino de datos y la unidad de control, sino también la resolución de conflictos asociados a dependencias de datos, riesgos de control y sincronización entre etapas.

En segundo lugar, el trabajo busca dotar al sistema de mecanismos de observabilidad y depuración que permitan inspeccionar su estado interno durante la ejecución. Para ello, se incorpora una \textit{Debug Unit} capaz de interactuar con el procesador mediante comandos enviados por UART, posibilitando tareas como la carga de programas en memoria, la ejecución paso a paso, la detención controlada, y el volcado de registros, memoria y latches del pipeline.

Finalmente, se busca validar el funcionamiento integral del sistema tanto en simulación como en FPGA, verificando que la ejecución de programas, el manejo de instrucciones especiales como \texttt{HALT}, y la interacción con la unidad de depuración se comporten de acuerdo con lo esperado. De esta forma, el trabajo no solo apunta a obtener un diseño funcional, sino también a comprender en profundidad la dinámica interna de un procesador segmentado y las herramientas necesarias para analizarlo y depurarlo.

\subsection{Descripción General del Sistema}

El sistema desarrollado consiste en un procesador basado en la arquitectura RISC-V de 32 bits, implementado como un \textit{pipeline} de cinco etapas clásicas: \textit{Instruction Fetch} (IF), \textit{Instruction Decode} (ID), \textit{Execute} (EX), \textit{Memory Access} (MEM) y \textit{Write Back} (WB). Cada etapa cumple una función específica dentro del ciclo de ejecución de instrucciones y se encuentra separada por registros intermedios que almacenan tanto datos como señales de control.

El procesador cuenta con memorias de instrucciones y datos, un banco de registros, lógica de forwarding y detección de hazards, y soporte para instrucciones aritméticas, lógicas, de acceso a memoria y de control de flujo dentro del subconjunto implementado. Además, incorpora señales de validez en los registros del pipeline, lo que permite distinguir entre instrucciones reales y burbujas introducidas por stalls o flushes, facilitando así un control más preciso de la ejecución.

Para permitir la interacción externa con el sistema, se utiliza una interfaz UART como canal de comunicación serie entre la FPGA y una computadora. A través de este enlace es posible enviar comandos y recibir respuestas desde un software de apoyo o una interfaz de usuario, sin necesidad de modificar directamente el hardware en cada prueba.

Sobre esta interfaz se apoya la \textit{Debug Unit}, módulo encargado de coordinar la programación de la memoria de instrucciones, el control de ejecución del procesador y la recolección de información interna del sistema. Esta unidad interpreta los comandos recibidos por UART y genera las señales necesarias para iniciar o detener la ejecución, avanzar un ciclo, reiniciar el procesador o extraer información de depuración. Gracias a esta integración, el sistema no solo ejecuta programas, sino que también puede ser observado y analizado de manera detallada, lo cual resulta especialmente útil durante las etapas de prueba, validación y demostración del funcionamiento del procesador.